\section{Introduction}

The companion write-up in this repository proved that the stationary equilibrium of
\citet{Aiyagari1994} is unique at his logarithmic calibration, and located exactly what stands in
the way for relative risk aversion $\mu\in\{3,5\}$. Every obstacle it met traces to the fixed point
of the infinite-horizon problem: a condition on the response of consumption to the interest rate
fails in a tail of wealth that the stationary distribution reaches; an induction that would prove
the response bound has a step multiplier $1/\Thorn>1$, so it cannot close as a supremum argument;
and the only floor on capital supply comes from an ergodic identity that loses a factor
$\|u'\|^2/\operatorname{Var}$. Overlapping generations remove all three at once. A household
that lives $J$ periods solves a finite backward induction from an explicit terminal condition, so
a multiplier $1/\Thorn$ per step compounds to $\Thorn^{-J}$, which at Aiyagari's numbers is about
$1.03$. Wealth is bounded by $J$ periods of accumulation, so any condition needs checking on a
known bounded region. And the distribution of the youngest cohort is common to every interest
rate, so an ordering of policies pushes forward to an ordering of the whole age profile of
distributions by induction on age, with no Doeblin atom and no uniqueness of a stationary
distribution.

The target of the programme is uniqueness of the stationary equilibrium of a stochastic
life-cycle economy in the manner of \citet{Huggett1996} at $\mu\in\{3,5\}$. Section~\ref{sec:lc}
records the first stage, the life-cycle household, which is formalised. Section~\ref{sec:prog}
records the rest of the programme, the classical results of \citet{Diamond1965} included, and
what each stage needs. As in the companion, every result stated with a formal name has been
checked by the Lean kernel against Mathlib; the development is at
\url{https://github.com/LeanEconomics/LeanEconomics}, directory \texttt{LeanEconomics/OLG}.
